\section{How the finite inequalities imply bounds for arbitrary laws}
\label{ext:finite-soundness}

The finite calculations bound functions on continuous parameter domains.
The first task in each application is to prove that the parameters of the
same extremal law belong to that domain. The calculations then verify
pointwise majorants or exclude all possible parameter values. They do not
replace the marginal by a finitely supported law.

\begin{lemma}[Tensor expectation for a common marginal]
\label{ext:tensor-soundness}
Let $K$ contain the support of a standardized law $P$, and let
$f_0=1,f_1,\ldots,f_d$ be integrable under $P$. Put $m_j=\E f_j(X)$.
For a finite tensor sum
\[
 T(x_1,\ldots,x_n)=\sum_\alpha t_\alpha
                       \prod_{i=1}^n f_{\alpha_i}(x_i),
 \qquad M_T=\sum_\alpha t_\alpha\prod_{i=1}^n m_{\alpha_i},
\]
one has $\E T(X_1,\ldots,X_n)=M_T$. If
\[
 T\ge\ind_{\{x_1+\cdots+x_n\le H\}}\quad\hbox{on }K^n,
 \qquad M_T-\Phi(H/\sqrt n)\le b\rho/\sqrt n,
\]
then the positive closed-CDF score of $P$ at $H$ is at most $b$.
\end{lemma}
\begin{proof}
Independence makes the absolute expectation of each product equal to the
product of its finite absolute expectations. Fubini's theorem and linearity
therefore evaluate the finite sum. Integrate the pointwise inequality and
multiply by $\sqrt n/\rho$, which is positive.
\end{proof}

A product with a centered factor has expectation zero if each factor
uses a different summand. Adding such products can improve a majorant
without changing $M_T$. The same marginal supplies every $m_j$;
one cannot choose different moment vectors for different factors.
Polynomial features of degree above three are used only after a support
bound has supplied their integrability.

\begin{lemma}[Exact positivity on a box]\label{ext:bernstein}
Let $p(t)=\sum_\alpha a_\alpha t^\alpha$ have coordinate degrees at most
$d_i$. On $[0,1]^r$ its tensor Bernstein coefficients are
\[
 b_\nu=\sum_{\alpha\le\nu}a_\alpha
                 \prod_{i=1}^r\frac{\binom{\nu_i}{\alpha_i}}
                                      {\binom{d_i}{\alpha_i}}.
\]
If all $b_\nu\ge0$, then $p\ge0$ throughout the closed box.
\end{lemma}
\begin{proof}
The identity
\[
 t^k=\sum_{j=k}^d\frac{\binom jk}{\binom dk}
              \binom dj t^j(1-t)^{d-j}
\]
follows from the binomial theorem after extracting $t^k$.
Apply it in each coordinate. The resulting basis functions are
nonnegative and sum to one, so $p$ is a convex combination of the
$b_\nu$ at each point of the box.
\end{proof}

Rational affine changes map general rational boxes to the unit box,
and exact subdivision can strengthen the coefficient test. The
polynomial identity, coordinate degrees and every retained child box
are part of the check. For a closed event the majorant must satisfy
$T-1\ge0$ on the event plane as well as its lower side. Checking
only $T\ge0$ on that plane would lose threshold atoms.

\begin{proposition}[Closed-cover soundness]\label{ext:closed-cover}
Suppose every law and threshold under consideration has a proved map
into a domain covered by finitely many closed root cells. At each root
there is a finite subdivision tree, and every split retains closed
children whose union contains the parent. Suppose that on every leaf
either a proved necessary condition excludes all actual mapped points,
or a valid whole-leaf inequality bounds their positive scores by $b$.
Then all actual scores in the root domain are at most $b$.
If instead every leaf excludes an assumed extremal score $C\ge b$,
no such extremizer belongs to the domain.
\end{proposition}
\begin{proof}
Every point of a parent belongs to a child, including boundary points.
Following children down a finite tree reaches a leaf. Its conclusion
therefore applies to the point. The root cover proves the assertion
on the whole domain.
\end{proof}

Each application must specify its support and moment bounds, its root
domain, and the inequality checked at a leaf. An empty-cell test must
exclude actual laws, not merely inconvenient numerical parameters.
Transcendental quantities require proved outward enclosures; divisions
require a positive denominator on the entire cell. Missing roots,
unresolved leaves and omitted endpoint cases prevent acceptance.
